%==============================================================================
%  v9 / sections/related.tex                                      [owner: A3]
%
%  DELIVERABLE STATUS: POINTS OUTLINE, not finished prose.
%
%  Coverage is the union of (a) v8:85-87's own TODO -- spectral phylogenetics,
%  low-rank matrix completion, clustering margins / range restriction,
%  macro-evolutionary tree-shape models -- and (b) the two bodies of work the
%  paper actually leans on but never surveys: community detection / NMI, and
%  entry-wise eigenvector analysis.
%
%  Each bullet ends with the same move: what the line gives us, and what it does
%  not cover -- so the section reads as a gap argument rather than a booklist.
%==============================================================================

\section{Related Work}\label{sec:related}

%------------------------------------------------------------------------------
\subsection*{Spectral phylogenetics and latent tree recovery}
%------------------------------------------------------------------------------
\begin{itemize}[leftmargin=1.4em, itemsep=3pt, topsep=2pt]
  \item \textbf{Direct ancestor of this work.} Spectral top-down recovery of
        latent tree models \cite{aizenbud2023spectral}: the log-det similarity
        matrix, the divide-and-conquer schema, and the clan-partition step we
        analyse. \emph{Gap:} the similarity matrix is assumed available in full,
        and imbalance is not a parameter of the guarantee.
  \item \textbf{Combinatorial substrate.} Splits, clans versus clades, and what a
        bipartition determines about a tree
        \cite{semple2003phylogenetics, wilkinson2007clades}.
  \item \textbf{Sequence-length requirements.} The ``few logs suffice'' line
        \cite{erdos1999logs} fixes how tree diameter and depth enter
        reconstructability; we reuse exactly those two quantities in
        \Cref{cor:tolerance}. \emph{Gap:} the budget accounted for is sequence
        length $\ell$, not the number of pairs evaluated.
  \item \textbf{Distance-based / numerical-taxonomy route.} Partition recovery
        from the leading eigenvector of a double-centred additive distance matrix
        \cite{Griffing2012}, resting on split decomposition of tree metrics
        \cite{BandeltDress1992}. \emph{Gap:} deterministic and fully observed; no
        sampling analysis. This is the line \Cref{thm:main-dist} extends.
\end{itemize}

%------------------------------------------------------------------------------
\subsection*{Low-rank matrix completion}
%------------------------------------------------------------------------------
\begin{itemize}[leftmargin=1.4em, itemsep=3pt, topsep=2pt]
  \item \textbf{The regime we sit in.} Exact recovery of a low-rank matrix from a
        uniformly random subset of entries, and the near-optimal
        $O(\mathrm{rank}\cdot m\log^2 m)$ sample complexities
        \cite{candes2008exact, candes2009power, recht2011simpler}; the noisy
        variant and inverse-probability reweighting
        \cite{candes2009matrix}.
  \item \textbf{Where coherence enters.} These results are governed by an
        incoherence assumption on the leading subspace, and coherent instances
        genuinely need more samples \cite{chen2014completing}. \emph{This is the
        conceptual borrowing of the paper:} imbalance is our coherence.
  \item \textbf{Gap.} The objective differs. Completion asks for the matrix in
        Frobenius or spectral norm; we ask only for the \emph{sign pattern of one
        eigenvector}, which is a weaker target but an $\ell_\infty$ one --- so
        neither the rates nor the norms transfer directly, and we do not run a
        completion algorithm at all.
\end{itemize}

%------------------------------------------------------------------------------
\subsection*{Clustering margins and range restriction}
%------------------------------------------------------------------------------
\begin{itemize}[leftmargin=1.4em, itemsep=3pt, topsep=2pt]
  \item \textbf{Noise thresholds for spectral clustering}
        \cite{Balakrishnan2011clustering}: the range-restricted similarity model
        and the margin between within- and between-cluster similarity that our
        $\rho$ (\eqref{eq:margin_definition}) is taken from.
  \item \textbf{Gap 1 --- balance.} The guarantees are stated with clusters of
        comparable size; balance sits in the analysis as a constant. Our $\eta$ is
        precisely that constant promoted to a variable and pushed to
        $\mathcal O(m)$.
  \item \textbf{Gap 2 --- observation model.} The similarity matrix is assumed
        fully observed; the noise is additive on entries, not missingness.
\end{itemize}

%------------------------------------------------------------------------------
\subsection*{Community detection and partition scoring}
%------------------------------------------------------------------------------
\begin{itemize}[leftmargin=1.4em, itemsep=3pt, topsep=2pt]
  \item \textbf{Two-block recovery as community detection.} Our problem is a
        planted two-community instance with a tree-generated similarity kernel;
        the spectral-clustering-under-blockmodels literature is the closest
        methodological neighbour \cite{fortunato2010community, qinrohe2013regularized}.
  \item \textbf{Degree correction is the same phenomenon under a different
        name.} Regularised spectral clustering under the degree-corrected
        blockmodel \cite{qinrohe2013regularized} addresses signal localising on
        low-degree nodes --- which is what coherence inflation does here; worth
        stating the correspondence explicitly.
  \item \textbf{Scoring.} NMI as the standard accuracy measure against a known
        planted partition, and its maximality characterisation
        \cite{strehlghosh2002cluster, fortunato2010community}, used in
        \Cref{cor:nmi}.
  \item \textbf{Gap.} Blockmodel guarantees are asymptotic in graph size with
        fixed community proportions; the imbalanced, sub-sampled, tree-generated
        case is not covered.
\end{itemize}

%------------------------------------------------------------------------------
\subsection*{Entry-wise eigenvector analysis}
%------------------------------------------------------------------------------
\begin{itemize}[leftmargin=1.4em, itemsep=3pt, topsep=2pt]
  \item \textbf{The baseline being moved away from.} Davis--Kahan's $\sin\Theta$
        theorem \cite{davis1970rotation}: the canonical eigenvector perturbation
        bound, and an $\ell_2$ one --- the reference point for the whole of
        \Cref{sec:main} and \Cref{app:compare}.
  \item \textbf{Beyond $\ell_2$.} Resolvent / Neumann-series expansions and
        $\ell_\infty$ eigenvector bounds
        \cite{eldridge2017unperturbed, fan2018eigenvector}, and the entry-wise
        analysis of low-expected-rank random matrices
        \cite{abbe2020entrywise}. \Cref{lem:neumann} is an instance of this
        machinery, specialised to the two-block Laplacian.
  \item \textbf{Concentration tool.} Matrix Bernstein
        \cite{tropp2012userfriendly}, used in \Cref{app:comp2} to control
        $\opnorm{\EL}$.
  \item \textbf{Gap.} These works target random-graph and covariance models; none
        instantiates the bound for a tree-generated similarity Laplacian, and none
        carries an imbalance parameter through to a sampling rate.
\end{itemize}

%------------------------------------------------------------------------------
\subsection*{Macro-evolutionary tree-shape models}
%------------------------------------------------------------------------------
\begin{itemize}[leftmargin=1.4em, itemsep=3pt, topsep=2pt]
  \item \textbf{The two null models.} The Yule / pure-birth process
        \cite{yule1925mathematical} and Kingman's coalescent
        \cite{kingman1982coalescent} --- the models our simulated trees are drawn
        from, and the source of the $\mathcal O(\log m)$ depth and diameter used
        in \Cref{cor:tolerance} and \Cref{asm:hbm}.
  \item \textbf{Their shape statistics.} Distributions of clade size and
        imbalance under Yule-type speciation \cite{steel2001properties}, and the
        Yule-to-today survey of tree-shape statistics \cite{aldous2001stochastic}
        --- i.e.\ the literature that tells us what $\eta$ actually looks like.
  \item \textbf{Why this matters for us, not just for biology.} Real trees are
        measurably \emph{more} imbalanced than either null model
        \cite{blum2006which}, so the imbalanced regime is the typical regime, not
        an adversarial corner. This is the empirical warrant for making $\eta$ the
        organising parameter.
\end{itemize}
